% Part: proof-theory
% Chapter: normalization
% Section: introduction

\documentclass[../../../include/open-logic-section]{subfiles}

\begin{document}

\olfileid{pt}{nor}{int}

\olsection{ভূমিকা}

ধরা যাক, আপনি শর্তসূত্র~$!A \lif !B$
!!{prove} করেছেন। \Elim{\lif} ব্যবহার করে
সেটি~$!B$-র !!a{proof}-এ কাজে লাগতে পারে।
অবশ্য এর জন্য~$!A$-র
!!a{proof}~$\delta_1$-ও লাগবে।
$!A \lif !B$-র আপনার !!{proof} সম্ভবত
\Intro{\lif} দিয়ে শেষ হয়: এতে
খোলা অনুমান~$!A$ থেকে~$!B$-র
!!a{proof}~$\delta_2$-কে
$!A \lif !B$-র প্রমাণে পরিণত করা হয়।
তাই হলে আপনার শেষ !!{proof}-এর রূপ
\[
\AxiomC{$\Discharge{!A}{x}$}
\RightLabel{$\delta_2$}
\DeduceC{$!B$}
\DischargeRule{\Intro{\lif}}{x}
\UnaryInfC{$!A \lif !B$}
\AxiomC{}
\RightLabel{$\delta_1$}
\DeduceC{$!A$}
\RightLabel{\Elim{\lif}}
\BinaryInfC{$!B$}
\DisplayProof
\]
এই কৌশলে আগের পাওয়া
$!A \lif !B$-র !!a{proof} আবার ব্যবহার
করায় কাজের সুবিধা হয়; কিন্তু শেষ
!!{proof}টি সরল নয়। সঙ্গে সঙ্গে নিষ্কাশনের
জন্যই $!A \lif !B$-কে প্রথমে ভূমিকা
করানো যেন ঘুরপথ। $!B$ !!{prove} করার
আরও সরাসরি পথ থাকার কথা।
% BN-SRC-574 TARGET-CORRECTION-BEGIN
\par\smallskip\noindent\textnormal{\small\textbf{উৎস-সংশোধন (BN-SRC-574):} উৎসগদ্যে $\delta_1$-কে $!A$ ও $!A\lif !B$ উভয়ের প্রমাণ বলা হয়েছে, অথচ প্রদর্শিত দ্বিতীয় শাখায় ভুলে $\delta_2$ বসেছে। $\delta_1$-কে $!A$-র এবং $\delta_2$-কে খোলা অনুমান থেকে $!B$-র প্রমাণ রেখে দ্বিতীয় শাখার চিহ্ন সংশোধিত।}
% BN-SRC-574 TARGET-CORRECTION-END

বাস্তবে এমন পথ সব সময়ই আছে। উপরের
মতো “ঘুরপথ”—!!a{introduction}
বিধির ঠিক পরে !!a{elimination}
বিধি—স্বাভাবিক নিষ্পাদনের !!{proof}s
থেকে বাদ দেওয়া যায়। একে
\emph{স্বাভাবিকীকরণ} বলে; ফল হলো
\emph{স্বাভাবিক} !!{proof} (অর্থাৎ
\emph{স্বাভাবিক রূপে} !!a{proof})।

সিকোয়েন্ট গণনায় কর্তন-নিষ্কাশন
উপপাদ্যের মতোই এই ফলের প্রমাণ কিছুটা
জটিল; অনেক ক্ষেত্র বিচার করতে হয়।
স্বজ্ঞাবাদী \Log{N1i}-র চেয়ে ধ্রুপদি
\Log{N1c}-র জন্য এটি প্রমাণ করা কঠিন।
সিকোয়েন্ট গণনায় কর্তন-নিষ্কাশন যেমন
দেখায় যে \Cut{} বিধি ব্যবহার করে
বাড়তি কিছু !!{prove} করা যায় না,
স্বাভাবিকীকরণও তেমনি দেখায় যে
ঘুরপথ-সহ যে ফল প্রমাণ করা যায়,
ঘুরপথ এড়িয়েও তা করা যায়। আরও মিল
আছে: একই ফলের ঘুরপথ-সহ সবচেয়ে
ছোট !!{proof}-এর চেয়ে স্বাভাবিক
!!{proof} অনেক দীর্ঘ হতে পারে,
যেমন সিকোয়েন্ট গণনায় কর্তনযুক্ত
!!{proof}s কর্তনহীন সবচেয়ে ছোট
!!{proof}-এর চেয়ে অনেক ছোট হতে পারে।
উপ-!!{formula} ধর্মের স্বজ্ঞাবাদী রূপে, স্বাভাবিক
!!{proof}-এ খোলা পূর্বধারণা ও সিদ্ধান্তের উপ-!!{formula}s,
পরিমাণসূচকের ক্ষেত্রে সেগুলির উপযুক্ত প্রতিস্থাপিত রূপ,
এবং প্রয়োজনে মিথ্যা-সূত্র ব্যবহার করেই
সিদ্ধান্ত !!{prove} করা যায়।
স্বাভাবিক নিষ্পাদনে এটি
স্বাভাবিকীকরণ থেকে আসে, যেমন সিকোয়েন্ট
গণনায় কর্তন-নিষ্কাশন থেকে আসে।
% BN-SRC-575 TARGET-CORRECTION-BEGIN
\par\smallskip\noindent\textnormal{\small\textbf{উৎস-সংশোধন (BN-SRC-575):} উৎসে উপসূত্র-ধর্মে শুধু সিদ্ধান্তের উপসূত্র বলা হয়েছে; পূর্বধারণা $P,P\lif Q$ থেকে $Q$ প্রমাণে $P$ ও $P\lif Q$ দরকার। তাই খোলা পূর্বধারণাগুলির উপসূত্রও শর্তে অন্তর্ভুক্ত।}
% BN-SRC-575 TARGET-CORRECTION-END
% BN-SRC-879 TARGET-CORRECTION-BEGIN
\par\smallskip\noindent\textnormal{\small\textbf{পাঠ-সংশোধন (BN-SRC-879):} খোলা পূর্বধারণা অন্তর্ভুক্ত করলেও আক্ষরিক উপসূত্রের
দাবি যথেষ্ট নয়: নঞর্থক নিষ্কাশন ও মিথ্যা থেকে সিদ্ধান্তে মিথ্যা-সূত্র
লাগতে পারে, আর পরিমাণসূচক বিধিতে প্রতিস্থাপিত রূপ আসে।
এই বক্তব্য স্বজ্ঞাবাদী রূপে সীমিত রাখা হয়েছে; ধ্রুপদি অব্যাহতি-বিধিতে
আক্ষরিক উপসূত্রের বাইরে নঞর্থক অনুমানও আসতে পারে।}
% BN-SRC-879 TARGET-CORRECTION-END

\end{document}
